\input qpd-bootstrap

\begin{problem}{20260808}
  \meta{name}{Binary Star System}
  \meta{setter}{Sophie Chen}
  \meta{score}{12}
  \meta{topic}{gravitation}
  \meta{difficulty}{4}
  \meta{tags}{gravitation, orbital motion, Kepler}

A binary star system consists of two stars orbiting their common centre of mass.  Star A has mass $M_A = 4\times10^{30}\,\mathrm{kg}$ and Star B has mass $M_B = 1\times10^{30}\,\mathrm{kg}$.  The straight-line distance between the two stars is $D = 5.0\times10^{12}\,\mathrm{m}$, and the gravitational constant is $G = 6.67\times10^{-11}\,\mathrm{N\,m^{2}\,kg^{-2}}$.  The two stars orbit their shared centre of mass in perfect uniform circular motion.

\begin{assumptions}
  \item The stars are point masses; only their mutual gravitational attraction acts.
  \item The orbits are circular and coplanar, about the common centre of mass.
\end{assumptions}

\begin{questions}
  \item \textbf{Equal and Opposite Forces.}
        Compare the magnitude of the gravitational pull that Star A exerts on Star B with the pull that Star B exerts on Star A.  State whether they are larger, smaller, or equal, and name the physics law that supports your conclusion.

  \item \textbf{The Zero-Force Point.}
        Along the line connecting Star A and Star B there is a unique point where the total gravitational force on a tiny test mass $m$ placed at that position cancels out to zero.
        \begin{enumerate}[label=\arabic*.]
          \item Let $x$ be the distance from Star A to this zero-force point.  Write the gravitational force-balance equation for the test mass.
          \item Solve numerically to calculate the value of $x$.
        \end{enumerate}

  \item \textbf{Orbital Period.}
        The zero-force point from part (b) is \emph{not} the system's centre of mass.  Both stars revolve around their common centre of mass with identical period $T$.
        \begin{enumerate}[label=\arabic*.]
          \item Write the equation equating the mutual gravitational force to the centripetal force for Star A.
          \item Write the matching force equation for Star B.
          \item Combine the two equations to derive the modified Kepler's Third Law for binary stars, then substitute the given numbers to solve for the orbital period $T$ of the system.
        \end{enumerate}

  \item \textbf{Orbital Speeds.}
        Derive and write a single equality relating $M_A$, $M_B$, $v_A$, and $v_B$.  No numerical calculation is needed here.
\end{questions}

\end{problem}

\begin{solution}{20260808}

\subsection*{Q1. Equal and Opposite Forces}

By Newton's Third Law the two pulls are equal and opposite: $\abs{\arr{F}_{A\to B}} = \abs{\arr{F}_{B\to A}}$.

\subsection*{Q2. The Zero-Force Point}

\textbf{(1)}  On the line AB the forces on $m$ are collinear and opposite, so they cancel when
\[
\frac{G M_A m}{x^{2}} = \frac{G M_B m}{(D-x)^{2}}
\;\Longrightarrow\;
\frac{M_A}{x^{2}} = \frac{M_B}{(D-x)^{2}}.
\]

\textbf{(2)}  Taking square roots and solving:
\[
x = D\,\frac{\sqrt{M_A}}{\sqrt{M_A}+\sqrt{M_B}}
  = (5.0\times10^{12})\,\frac{2\times10^{15}}{3\times10^{15}}
  = \boxed{3.3\times10^{12}\ \mathrm{m}}.
\]

\subsection*{Q3. Orbital Period}

Let $r_A + r_B = D$ be the orbital radii about the common centre of mass, with $\omega = 2\pi/T$.  Equating the mutual force to the centripetal force for each star,
\[
\frac{G M_A M_B}{D^{2}} = M_A\omega^{2} r_A = M_B\omega^{2} r_B.
\]
With $M_A r_A = M_B r_B$ and $r_A+r_B=D$, we get $r_A = \dfrac{M_B D}{M_A+M_B}$; substituting and using $\omega = 2\pi/T$ gives the modified Kepler's Third Law:
\[
\boxed{T^{2} = \frac{4\pi^{2}D^{3}}{G(M_A+M_B)}}.
\]
Numerically, with $M_A+M_B = 5.0\times10^{30}\,\mathrm{kg}$:
\[
T = 2\pi\sqrt{\frac{(5.0\times10^{12})^{3}}{(6.67\times10^{-11})(5.0\times10^{30})}}
  = \boxed{3.8\times10^{9}\ \mathrm{s}} \approx 122\ \text{yr}.
\]

\subsection*{Q4. Orbital Speeds}

Same angular speed means $v_A/v_B = r_A/r_B = M_B/M_A$, hence
\[
\boxed{M_A v_A = M_B v_B}.
\]

\end{solution}

\input qpd-end
